\documentclass[11pt,a4paper]{article}
\usepackage{fontspec}
\setmainfont{Noto Serif}
\setsansfont{Noto Sans}
\usepackage{babel}
\babelprovide[main,import]{russian}
\usepackage{amsmath,amssymb,mathtools,booktabs,geometry,fancyhdr,microtype,xcolor,hyperref}
\geometry{left=21mm,right=21mm,top=21mm,bottom=21mm}
\definecolor{ink}{HTML}{173D57}
\hypersetup{colorlinks=true,urlcolor=ink,linkcolor=ink,pdftitle={BGS-F. Второй порядок и простота собственных фаз}}
\pagestyle{fancy}\fancyhf{}\lhead{\small BGS-F}\rhead{\small Второй порядок и простота собственных фаз}\cfoot{\thepage}
\setlength{\headheight}{15pt}\setlength{\parindent}{1em}\setlength{\parskip}{2.5pt}
\setlength{\emergencystretch}{3em}\tolerance=3000
\newcommand{\Tr}{\operatorname{Tr}}
\newcommand{\EE}{\mathbb E}
\newcommand{\R}{\mathbb R}
\newcommand{\Z}{\mathbb Z}
\newcommand{\C}{\mathbb C}
\newcommand{\si}{\mathrm{Sine}_2}
\newcommand{\eps}{\varepsilon}
\newcommand{\sinc}{\operatorname{sinc}}
\newcommand{\sect}[1]{\section*{\sffamily\color{ink}#1}}
\newcommand{\subt}[1]{\par\medskip\noindent\textbf{#1.}\quad}
\newcommand{\T}{\mathbb T}
\newcommand{\dist}{\operatorname{dist}}
\newcommand{\Ad}{\operatorname{Ad}}
\newcommand{\zp}{\mathbb Z_7}
\newcommand{\QQ}{\mathbb Q}
\newcommand{\kk}{\overline{\mathbb Q}}
\newcommand{\ord}{\operatorname{ord}}
\newcommand{\diag}{\operatorname{diag}}
\begin{document}
\noindent{\small\sffamily\color{ink} МАТЕМАТИЧЕСКАЯ ФИЗИКА\hfill 22 НОЯБРЯ 2026}
\begin{center}{\LARGE\sffamily Второй порядок и простота\\[3pt]собственных фаз}\\[6pt]
{\large Точный фиксированный кик на размерностях \(7^k\)}
\end{center}
Рассматривается исходная матрица BGS-F
\[
    U_N=\exp(i10^{-80}N A_N)M_N,\qquad N=7^k.
\]
Пусть \(s_N\) есть число фаз кратности один. Для линейного \(M_N\) после снятия общей фазы узлы имеют период \(m=3\cdot7^{k-1}\). Узлы с примитивным \(7\)-адическим характером при \(k\ge2\) являются двойными.

\subt{Полное расщепление примитивных двойных уровней}
В аналитическом семействе \(e^{itA_N}M_N\) каждая из этих пар разделяется в первом или во втором порядке при \(t=0\), для каждого \(k\ge2\). В частности, при исходном фиксированном \(10^{-80}\) число простых фаз не меньше \(36N/49\) уже за счёт этих ветвей. Само число простых фаз относится к спектру полного унитарного оператора. Для перехода на точный параметр используется инъективность алгебраической специализации конечных экспоненциальных полиномов.

\subt{Оценка доли простых собственных фаз}
С учётом кластеров более высокой кратности, включая дополнительное расщепление на глубине два, получаем
\[
   \boxed{\displaystyle
   \liminf_{k\to\infty}\frac{s_{7^k}}{7^k}\ge
                    \frac{18867}{19208}>0.9822 .}
\]
Предыдущая граница \(345/392\) составляла примерно \(0.8801\). Точная оценка совпадений спектральных разностей с нулём даёт
\[
   \limsup_{k\to\infty}\lim_{T\to\infty}
        \frac1{T7^k}\sum_{r=1}^T|\Tr U_{7^k}^r|^2
             \le \frac{58663}{57624}<1.0181 .
\]

\subt{Механизм второго порядка}
При скалярном первом коэффициенте остаётся матрица на двойном уровне, содержащая резольвенту линейного пропагатора. После редукции по простому над \(2\) равенство \(\lambda+\mu=\lambda-\mu\) убирает знаменатель для каждой другой линейной фазы \(\mu\). Проверка нескалярности сводится к внедиагональному элементу квадрата целочисленно нормированного \(A_N\). Во всех семи нужных классах по модулю \(49\) один и тот же фазовый остаток \(11\) имеет единственный ненулевой вклад. Это даёт доказательство на всех \(k\ge4\) из одного конечного фазового сертификата. Модули \(49\) и \(343\) разобраны отдельно в конечных полях.

Исходная задача требует слабого предела локального спектрального процесса. Оценки кратностей определяют точные совпадения, а для предела нужно также контролировать ненулевые расстояния порядка \(1/N\). Порядок пределов во второй формуле указан явно.

\newpage
\sect{1. Дополнение Шура и редукция над двойкой}
Обозначим \(B=4000A_N\). Общее скалярное вращение \(M_N\) удаляем. После симплектической замены меток \((v_1,v_2)\mapsto(a,b)\) оператор \(M\) действует как дилатация с множителем \(\gamma=3+2\rho\) и \(M^m=I\), где
\[
 \rho^2=2,\quad \rho\equiv3\pmod7,qquad
  (a,b)=\left(v_1+\rho v_2,\frac{v_1-\rho v_2}{-2\rho}\right).
\]
Замена меток сохраняет симплектическую форму и реализуется точным унитарным сопряжением Вейлева базиса. Для \(\lambda^m=1\) положим \(P=m^{-1}\sum_{r=0}^{m-1}\lambda^{-r}M^r\), \(Q=I-P\). Аналитическое семейство \(V(s)=e^{sB}M\) изучаем при \(s\) около нуля. Исходной матрице соответствует \(s=iN/(4000\cdot10^{80})\).

Пусть \(G=PBP\) рассматривается на \(\operatorname{ran}P\). Для \(z=\lambda+O(s)\) обычное дополнение Шура матрицы \(V(s)-zI\) на \(P\), после деления на \(\lambda\), имеет вид
\[
  (1-z/\lambda)I+sG+\frac{s^2}{2}(G^2+J)+O(s^3),
\qquad
 J=PBQ\,\frac{\lambda+M}{\lambda-M}\,QBP.                 \tag{1}
\]
Дробь определена на \(\operatorname{ran}Q\). Действительно, диагональный блок квадратичного порядка даёт \(s^2PB^2P/2\), а смешанное произведение даёт \(s^2PBQM(\lambda-M)^{-1}QBP\). Тождество \(PB^2P=G^2+PBQBP\) даёт (1). У \(J\) косоэрмитова симметрия, так как \((\lambda+\mu)/(\lambda-\mu)\) чисто мнимо при \(\mu\ne\lambda\) на окружности. При \(\dim P=2\) нескалярный \(G\) разделяет обе ветви в первом порядке. Если \(G=cI\), достаточно нескалярного \(J\). Два корня при подстановке \(z/\lambda=1+sc+s^2\beta\) разделяет различие собственных чисел \((c^2I+J)/2\). Это следует из принципа аргумента в двух малых непересекающихся кругах по \(\beta\). Рассуждение применимо при малом мнимом \(s\), когда \(V(s)\) унитарна.

\subt{Локальное кольцо}
Работаем над \(\QQ(i,\zeta_3,\zeta_N)\) и в локализации его кольца целых по простому над \(2\). Поскольку \(m\) нечётно, знаменатель \(m\) обратим. Все различные \(m\)-е корни единицы остаются различными после редукции. Произведения их разностей в каждой строке равны производной \(X^m-1\) в данном корне и являются единицами. Поэтому выражение для \(J\) корректно редуцируется и удовлетворяет
\[
          \overline J=\overline{PB^2P}-\overline G^{\,2}.     \tag{2}
\]
Черта обозначает редукцию. Это равенство проверяется на каждом линейном проекторе внутри \(Q\), где дробь в (1) редуцируется в 1. Если \(G\) скалярен, его внедиагональный элемент равен нулю, как и внедиагональный элемент \(G^2\). Тогда ненулевой элемент \(\overline{PB^2P}_{+,-}\) свидетельствует о нескалярности \(J\) в характеристике ноль. Фактор \(1/2\) в (1) вынесен за скобки и в редукции не используется.

\newpage
\sect{2. Единый фазовый свидетель по модулю 49}
Пусть \(\mathcal L(x)=\log(x^6)/6\) на \(\Z_7^\times\), \(\omega(3)=\zeta_3\) для характера по модулю 7. Базис узла \(\lambda\) при \(7\nmid\ell\) дан двумя векторами
\[
   \omega(x)^h e_{7^k}(\ell\mathcal L(x))
\]
на квадратной и неквадратной орбите единиц соответственно. Здесь \(e_n(y)=e^{2\pi i y/n}\), \(h\bmod3\), \(\ell\bmod7^{k-1}\). Векторы оставлены без нормировки, а коэффициент сжатия вычисляется с делением на размер орбиты \(m\). Это тот же матричный элемент, что при нормировке обоих векторов, и все числа редуцируются над двойкой.

В \(B\) коэффициенты исходных меток являются гауссовыми целыми. В редукции \(i\mapsto1\), и остаётся 12 меток с \(\Re(c)+\Im(c)\) нечётным, где \(c\) есть коэффициент \(B\). Произведение двух выживших меток \(v,w\) даёт результирующую метку \(v+w\) и фазу \(e_N(\overline2\sigma(v,w))\), \(\sigma(v,w)=v_1w_2-v_2w_1\). Сокращаем совпадающие тройки \((v+w,\sigma(v,w))\) по чётности. Получается 108 троек. Все исходные компоненты кокцикла имеют модуль не больше 18, как и компоненты результирующей метки. Поэтому равенство троек по модулю \(49\) совпадает с целым равенством. Список выводится из исходного символа при каждом запуске сертификата.

Для данной тройки \((v,c_0)\) и \(y=x+a\) критическое значение во внедиагональном элементе ищется среди \(x\) с квадратным \(x\bmod7\) и неквадратным \(y\bmod7\):
\[
   bxy=\ell a,\qquad
   \Phi(x)=b(x+\overline2 a)+\overline2 c_0
                          +\ell[\mathcal L(y)-\mathcal L(x)] .  \tag{3}
\]
При \(k\ge4\) необходимое сравнение для стационарного вклада выполняется по модулю \(49\). Это следует разбиением суммы на классы по \(7^{\lceil k/2\rceil}\) и точным линейным разложением фазы на каждом таком классе по модулю \(7^k\). Если \(a\equiv0\bmod7\), переход между противоположными орбитами невозможен. Иначе можно перечислять \(x\bmod49\) и вычислять \(\ell=bxy/a\) и \(\Phi(x)\bmod49\). Для каждого указанного остатка \(\ell\) фазу 11 даёт ровно одна тройка вместе с одним корнем. Во всех столбцах метка \(v=(0,10)\), а кокцикл \(c_0=0\).
\begin{center}
\begin{tabular}{lrrrrrrr}
\toprule
\(\ell\bmod49\) & 4 & 11 & 18 & 25 & 32 & 39 & 46\\
\(x\bmod49\) & 39 & 4 & 18 & 32 & 46 & 11 & 25\\
\(\Phi(x)\bmod49\) & 11 & 11 & 11 & 11 & 11 & 11 & 11\\
\bottomrule
\end{tabular}
\end{center}
Каждый найденный корень прост, поскольку \(b(2x+a)\) обратимо по модулю 7. При точном подъёме такой корень даёт экспоненту критического значения, умноженную на степень семи, конечный характер и возможную квадратичную сумму Гаусса по модулю 7. Последняя обратима над 2, поскольку её произведение с сопряжением равно 7. Поэтому коэффициент отдельного вклада после редукции не обнуляется. На следующей странице доказано, почему выделение фазы 11 исключает сокращение с остальными вкладами.

\newpage
\sect{3. Циклотомическая независимость и малые модули}
\subt{Независимость фазовых остатков}
В характеристике два \(\ord_{7^k}(2)=3\cdot7^{k-1}\). Порядок по модулю 7 равен 3 и \(2^3-1=7\), так что утверждение следует биномиальным поднятием. Присоединение \(\overline\zeta_3\) добавляет множитель 2 к степени конечного поля. Следовательно, при \(k\ge4\)
\[
  [\mathbb F_2(\overline\zeta_N,\overline\zeta_3):
       \mathbb F_2(\overline\zeta_N^{49},\overline\zeta_3)]=49 .
\]
Степени \(\overline\zeta_N^j\), \(0\le j<49\), являются базисом над меньшим полем. Вклад с фазой 11 по модулю 49 лежит в одной базисной компоненте, а остальные стационарные вклады в других. Для этого достаточно группировать фазы до точного подъёма, потому что по (3) фаза по модулю 49 постоянна на каждом классе \(x\bmod49\). Элементы группы 11 при точном поднятии дают один ненулевой ведущий коэффициент, включающий гауссов множитель из меньшего поля. Следовательно, \(\overline{PB^2P}_{+,-}\ne0\) при всех \(\ell\equiv4\bmod7\), всех \(h\) и всех \(k\ge4\). Для остальных пяти примитивных классов первый порядок уже даёт нескалярный \(G\) по точному сертификату по модулю 343 в части о кратностях при фиксированном кике.

\subt{Модули 49 и 343}
Для \(k=2,3\) используем точное суммирование исходных фаз в конечном поле. Положим \(s_0=7^{k-1}\) и возьмём модель
\[
  \mathbb F_4[X]/(X^{3s_0}+X^{s_0}+1),\qquad
  \mathbb F_4=\mathbb F_2[\omega]/(\omega^2+\omega+1).
\]
Многочлен является одним из множителей \(\Phi_{7^k}\), так как \(Y^3+Y+1\) делит \(\Phi_7(Y)\) по модулю 2. Его корни имеют порядок \(7^k\). Порядок 4 по модулю \(7^k\) равен \(3s_0\), так что этот множитель неприводим над \(\mathbb F_4\). Каждый вклад с показателем \(L+q s_0\), \(0\le L<s_0\), сводится к трём битам по остатку \(Y^q\bmod(Y^3+Y+1)\) с возможным весом \(\omega^h\) после учёта характерного отношения. Данные по \(x\) перечисляются прямо, стационарность здесь не требуется.

При \(k=2\) первый порядок для проверки требует дополнения на \(\ell=4,6\). Подстановка показывает, что модульный внедиагональный элемент \(PB^2P\) ненулевой для каждого из этих \(\ell\) и трёх \(h\). При \(k=3\) достаточно проверить семь остатков \(\ell\equiv4\bmod7\) и все \(h\); они дают ненулевой элемент. Для остальных \(\ell\) и \(h\) сохраняется свидетель первого порядка, вычисляемый над циклотомическим полем характеристики ноль. Итак, все \(18\) примитивных двойных узлов при \(N=49\), все \(126\) при \(N=343\) и все примитивные двойные узлы при \(k\ge4\) расщепляются при малом ненулевом \(t\).

\subt{Число ветвей}
На единицах есть две орбиты, каждая длины \(m=3\cdot7^{k-1}\). Примитивных характеров \(3(7^{k-1}-7^{k-2})=18N/49\). Их линейные фазы не встречаются вне единичного слоя. После расщепления удвоенное число узлов даёт \(36N/49\) простых ветвей.

\newpage
\sect{4. Дополнительное расщепление на глубине два}
Запишем \(\ell=7^d\mu\) при \(7\nmid\mu\). Для фиксированного \(d\ge1\) и больших \(k\) кластер содержит \(2(d+1)\) ветвей, по две на глубинах \(0,\ldots,d\), где глубина равна \(\nu_7(x)\). Первая производная в исходных углах равна сжатию \(A_N\). Арифметический расчёт в части о кратностях при фиксированном кике даёт четыре различных ненулевых значения на основном блоке глубин \(0,d\). Они отделены от остальных блоков. При \(d=1\) этим исчерпываются все ветви. При \(d=2\) остаётся двухмерный блок \(D\) на глубине 1. Здесь уточним, когда он нескалярен.

Вклад в \(D\) дают только \(v=(0,\pm7)\), с равными вещественными коэффициентами \(9/20\). После \(x=7z\), \(a=7a_0\), \(b=7b_0\) сокращается общий множитель \(7^2\) из фазы. Получаем те же две метки с разделёнными координатами для \((0,\pm1)\), параметр характера \(\mu\) и проводник \(7^{k-2}\). При усреднении по \(z\bmod7^{k-1}\) каждый класс по модулю проводника повторяется семь раз. Это влияет только на общий ненулевой множитель. Оператор коммутирует с отражением, которое меняет местами две орбиты, так что \(D=\alpha I+\beta R\), где \(R\) переставляет орбиты. Достаточно доказать \(\beta\ne0\). Для квадратного \(z\) и неквадратного \(z+a_0\) стационарные классы таковы.
\begin{center}
\begin{tabular}{cccc}
\toprule
Метка после деления & \((z,z+a_0)\bmod7\) & \(\mu\bmod7\) & Корень\\
\midrule
\((0,-1)\) & \((1,5),(2,6)\) & 2 & простые\\
\((0,1)\) & \((2,5)\) & 4 & двойной\\
\bottomrule
\end{tabular}
\end{center}
При \(\mu\equiv2\) две фазы различны по модулю 7: остатки \(2,5\). Поэтому при больших \(k\) оба вклада не сокращаются в циклотомическом расширении степени 7 по сравнению с предыдущим проводником.

При \(\mu\equiv4\) у единственной метки кладём \(y=z+a_0/2\). Уравнение принимает вид \(y^2=\Delta\), где
\[
         \Delta=a_0^2/4+\mu a_0/b_0 .
\]
В этом классе \(\Delta\) равномерна в \(7\Z_7\). При \(\nu_7(\Delta)=2j\ge2\) и квадратном ведущем разряде есть два стационарных корня \(\pm z_*\) с \(\nu_7(z_*)=j\). Разность критических фаз имеет порядок \(3j\). Действительно, производная фазы равна \(b_0(y^2-z_*^2)/(y^2-a_0^2/4)\), и ведущая разность первообразных в двух корнях равна \(-4c z_*^3/3\) с \(c=-4b_0/a_0^2\) обратимым. Для следующих членов порядок строго больше \(3j\) из оценки \(7^{2r}>2r+3\) при \(r\ge1\). Для \(k-2\ge4j+2\) точная стационарная фаза даёт два ненулевых вклада в разных циклотомических классах по модулю \(7^{3j+1}\). Значит \(\beta\ne0\). Остальные \(\mu\equiv4\) с нечётной оценкой \(\Delta\) или неквадратным ведущим разрядом имеют нулевой внедиагональный элемент при фиксированном порядке \(\Delta\) и достаточно большом проводнике, поскольку стационарного корня нет. Усечение высоких порядков имеет стремящуюся к нулю долю. Условная доля квадратных \(\Delta\) равна
\[
           \sum_{j\ge1}\frac{1-1/7}{2}\,7^{-(2j-1)}=\frac1{16}.
\]
Итого \(D\) добавляет две простые ветви в доле \((1+1/16)/6=17/96\) кластеров глубины два.

\newpage
\sect{5. Исходный кик и точные совпадения фаз}
\subt{Перенос на силу \(10^{-80}\)}
Матрицы \(A_N,M_N\) алгебраические. Элементы \(e^{itA_N}M_N\) и его обратной матрицы принадлежат кольцу конечных экспоненциальных сумм с алгебраическими частотами и коэффициентами. Это следует спектральным разложением \(A_N\). По теореме Линдемана--Вейерштрасса специализация в любой ненулевой вещественной алгебраической точке \(t\) сохраняет ненулевые элементы кольца. Поэтому сохраняется максимальный ранг каждой матрицы над ним. Для числа различных корней подходит матрица \((\Tr U(t)^{a+b})_{a,b=0}^{N-1}\). Она имеет ранг, равный числу различных корней, по факторизации Вандермонда. В точках максимального ранга кратности локально постоянны как мультимножество. Дополнение к дискретному исключительному множеству связно в \(\mathbb C\). Следовательно, все точки общего ранга наследуют по меньшей мере все простые ветви, полученные вблизи нуля. Это применимо к \(t=N/10^{80}\) без ограничения его роста. Для суммы квадратов кратностей отдельно достаточно переноса ранга \(\Ad(U(t))-I\), поскольку её ядро есть коммутант.

\subt{Доля простых фаз}
Число кластеров глубины \(d\) на единицу размерности равно \(18/7^{d+2}\) для каждого фиксированного \(d\). Складываем \(36/49\) из примитивных двойных узлов, четыре простые ветви на каждом кластере положительной глубины и дополнительную долю на глубине два. Суммирование сначала ведём по конечному набору глубин, потом переходим к \(k\to\infty\) и увеличиваем отсечку. Получаем
\[
 \liminf\frac{s_{7^k}}{7^k}\ge
 \frac{36}{49}
   +4\frac{18}{49}\sum_{d\ge1}7^{-d}
   +2\frac{18}{7^4}\frac{17}{96}
 =\frac{48}{49}+\frac{51}{19208}=\frac{18867}{19208}.
\]

\subt{Длительное среднее}
Для каждой фиксированной унитарной матрицы среднее \(T^{-1}\sum_{r=1}^T|\Tr U^r|^2\) стремится к сумме квадратов кратностей. Это геометрическая сумма по отношению фаз. Для линейного \(M_N\) эта сумма равна \((8N-5)/3\). Расщепление каждого примитивного двойного узла уменьшает сумму на 2. Для кластеров \(d\ge1\) гарантированные уменьшения из расчёта первых производных равны \(12\), \(28\), \(32d-44\) при \(d=1\), \(d=2\), \(d\ge3\) соответственно. Они получаются из четырёх простых ненулевых ветвей и остаточных блоков. На глубине два в только что посчитанной доле дополнительно вычитается 2. Следовательно,
\[
 \limsup_{k\to\infty}\lim_{T\to\infty}
 \frac1{T N}\sum_{r=1}^T|\Tr U_N^r|^2
 \le\frac83-\frac{36}{49}
   -\frac{18}{49}\left(\frac{12}7+\frac{28}{49}
                +\sum_{d\ge3}\frac{32d-44}{7^d}\right)
   -\frac{51}{19208}
 =\frac{58663}{57624}.
\]
Граница относится к фиксированной исходной силе, но время усреднения сначала стремится к бесконечности при фиксированном \(N\).

\newpage
\sect{6. Воспроизведение и связь с BGS-F}
\subt{Целочисленная часть}
Программа \texttt{exact\_splitting/check.py} читает вейлевы коэффициенты из исходной формулы. Она вычисляет квадрат над двойкой, находит фазовые остатки по модулю 49, проверяет простоту выбранных стационарных корней и выписывает ненулевые коэффициенты при малых проводниках. Отдельно проверяются таблица для блока глубины два, рациональные константы и сертификаты первого порядка при \(N=49,343\). Архив содержит точный вывод в \texttt{exact\_splitting/data/certificates.json}. Режим \texttt{--verify} пересчитывает его и проверяет побайтовое совпадение.

\subt{Независимая матричная часть}
Программа \texttt{exact\_splitting/check\_matrix.py} проверяет квадрат оператора прямым двухшаговым суммированием в координатном базисе над конечным полем при \(N=49\). При \(N=49,343\) дополнительно строится полный \(e^{itA_N}M_N\) в малых вспомогательных параметрах для проверки матрицы (1). Сверяются косоэрмитовость квадратичного коэффициента и порядок остатка при делении по параметру. В точном доказательстве исходная сила переносится алгебраически. Численный тест не используется для вывода об исходном пределе.

\subt{Где проходит граница доказанного}
Нулевые расстояния в спектре полного оператора теперь контролируются на большой доле фаз при всех \(k\). Локальный парный тест зависит от расстояний, сравнимых с \(1/N\), и от обычных следов при временах \(r\asymp N\). Простота собственных значений не предотвращает скопления различных фаз на таких расстояниях. Теорема Линдемана--Вейерштрасса сохраняет точный ненулевой ранг и не даёт оценки расстояний требуемого масштаба. В частности, оценки этой части не доказывают сходимость к \(\si\) или её нарушение при фиксированном \(10^{-80}\).

\subt{Проверяемые исходные утверждения}
Для переноса использована теорема Линдемана--Вейерштрасса о линейной независимости экспонент различных алгебраических чисел над алгебраическим замыканием \(\QQ\). В отчёте также используются явные расчёты четырёх внешних ветвей и отделения блоков, выписанные в разделах 7--10 предшествующей части о кратностях. Для них точные сертификаты поставляются в \texttt{spectrum/data/checks.json}.

\medskip
\noindent
\href{https://encyclopediaofmath.org/wiki/Lindemann_theorem}{Lindemann theorem, Encyclopedia of Mathematics.}

\end{document}
